\documentclass[11pt,a4paper,twocolumn]{article}

% Page layout
\setlength{\parindent}{.5cm}
\setlength{\columnsep}{.5cm}
\usepackage[left=2cm,right=2cm,top=2cm,bottom=2cm]{geometry}

% Encodings, characters, and typesetting
\usepackage[utf8]{inputenc}
\usepackage[T1]{fontenc}
\usepackage[english]{babel}
\usepackage{csquotes}
\usepackage{newtxtext,newtxmath} % access Times New Roman
\usepackage{microtype} % prettifies typesetting and prevents some overfulls

% Figures, tables, and captions
\usepackage{graphicx}
\usepackage{booktabs}
\usepackage[font=small,labelfont=bf,labelsep=period,singlelinecheck=false]{caption}
\graphicspath{{graphics/}}

% Title page
\usepackage{authblk}
\renewcommand\Affilfont{\small}

% References
\usepackage[backend=biber,style=nature]{biblatex}
\addbibresource{references.bib}

% Ignore underfull warning from page transition 1 to 2
\vbadness=10000

\title{\bf Social media reveals population dynamics of dysphoric dreaming}

\author[1,2,3]{Remington Mallett\thanks{correspondence: remington.mallett@u.northwestern.edu}}
\author[4]{Laura Sowin}
\author[2,3]{Michelle Carr}
\affil[1]{Department of Psychology, Northwestern University, USA}
\affil[2]{Department of Psychiatry and Addictology, University of Montreal, CAN}
\affil[3]{Center for Advanced Research in Sleep Medicine, Integrated University Health and Social Services Centres, CAN}
\affil[4]{Media Lab, Massachusetts Institute of Technology, USA}

\date{}

\begin{document}

\maketitle

\section*{Abstract}

Nightmares disrupt sleep health and predict future psychiatric diagnosis, yet reliable population estimates and their fluctuations over time are difficult to obtain. Here, we observe an increase in dysphoric dreaming shared on Reddit immediately following the World Health Organization's declaration of COVID-19 as a global pandemic. This ``digital dream surveillance'' approach might offer the field of sleep medicine a low-cost and real-time system for monitoring population sleep health.

\bigskip
\noindent {\bf Keywords:} nightmares, sleep, natural language processing, digital health surveillance, Reddit

\bigskip
\noindent
\textbf{Resource availability.} All code used to extract and analyze data is available in a public GitHub repository (https://github.com/remrama/covidreams). All data used in the current study is available in a public Zenodo repository (https://zenodo.org/records/18940545).

\bigskip
\noindent
\textbf{Author contributions.}
RM and MC designed the study. RM analyzed the data. RM and LS collected and curated the data. RM and MC wrote the manuscript.

\bigskip
\noindent
\textbf{Competing Interests.}
The authors declare no potential competing interests.

\bigskip
\noindent
\textbf{Acknowledgements.}
RM would like to thank Ashwini Ashokkumar for invaluable conversation and feedback related to the current project.

\newpage
\section*{Introduction}

Nightmares and other bad dreams incur major sleep health costs. Nightmare recallers suffer reduced sleep efficiency~\supercite{simor2012disturbed}, depressed mood~\supercite{blagrove2004relationship}, increased suicide risk~\supercite{que2022nightmares}, and are more likely to receive further psychiatric diagnoses than their counterparts~\supercite{wong2023prevalence}. Despite their clinical significance, reliable and ecological estimates of dysphoric dreaming in the population are difficult to obtain due to variable definitions~\supercite{worley2021epidemiology}. Thus, there is high demand in sleep medicine for a naturalistic approach to measuring population-level dysphoric dreaming over time and with multiple definitions.

\par
Here, we tested the ability to detect dynamic changes in multiple classes of dysphoric dreaming~\supercite{siclari2020dreams} -- nightmares and anxious dreaming -- on social media using natural language processing~\supercite{elce2021language}. Prior survey studies suggest that dysphoric dreaming increased during the first wave of the COVID-19 pandemic~\supercite{solomonova2021stuck}, but such studies have limited temporal precision and low ecological validity. Utilizing natural and unprompted human discourse from social media offers many advantages over retrospective survey studies~\supercite{lazer2021meaningful}, including reduced reporting bias and real-time monitoring that allows for preventative rather than reactive medicine. This digital health surveillance approach has been used frequently to track population health in psychiatry and other fields~\supercite{shakeri2021digital}, but is only beginning to gain traction in sleep health~\supercite{mciver2015characterizing,leypunskiy2018geographically}.

\begin{figure*}[t]
    \centering
    \includegraphics{figure.pdf}
    \caption{\textbf{Monitoring dysphoric dreaming on Reddit surrounding the declaration of COVID-19 as a global pandemic.} (A) Dream reports were extracted from Reddit by scraping \textit{r/Dreams} and filtering out posts that did not indicate a dream with post ``flair,'' leaving only dream posts. Anxious dreaming was quantified as the percentage of anxious words from the body text of dream posts (i.e., dream reports), and nightmares were identified as dream posts whose titles included a nonzero percentage of nightmare words. (B) Only dream-flaired posts surrounding the WHO's declaration of COVID-19 as a global pandemic were included in analysis. (C) Anxious dreaming -- the percentage of LIWC anxiety words -- increased following the WHO's declaration of COVID-19 as a global pandemic. Data line is smoothed over a 7-day moving average for visualization purposes only. Shaded error represents SEM. (D) More nightmares were posted to \textit{r/Dreams} in the month following the WHO's declaration of COVID-19 as a global pandemic than the month preceding. Error bars represent bootstrapped ($N = 2,000$) $95\%$ confidence intervals. (E) After the WHO's declaration of COVID-19 as a global pandemic, the frequency of COVID-related news posted on Reddit was predictive of anxious dreaming the following week.}
    \label{fig:summary}
\end{figure*}

\begin{table*}[t]
    \caption{\textbf{Statistics for main and control analyses.} Our primary analyses included only \textit{r/Dreams} posts that were flaired as dream reports surrounding March 11\textsuperscript{th}, 2020. To control for potential effects of seasonality, we ran the same analyses using only \textit{r/Dreams} posts that were flaired as dream reports surrounding March 11\textsuperscript{th}, 2019. To control for potential changes in common language, we ran the same analyses using only \textit{r/Dreams} posts that were not flaired as dream reports surrounding March 11\textsuperscript{th}, 2020. See Methods for descriptions of $B_1$, $B_2$, and $B_3$. Bolded values are statistically significant at $p < .05$.}
    \label{tab:statistics}
    \small
    \begin{tabular}{p{0.17\textwidth} p{0.04\textwidth} p{0.06\textwidth} p{0.04\textwidth} p{0.06\textwidth} p{0.04\textwidth} p{0.07\textwidth} p{0.04\textwidth} p{0.07\textwidth} p{0.04\textwidth} p{0.04\textwidth}}
        \toprule
        {} & \multicolumn{6}{l}{Anxiety regression} & \multicolumn{2}{l}{Anxiety corr.} & \multicolumn{2}{l}{Nightmares chi-sq.} \\
        Text source & $B_1$ & $p$ & $B_2$ & $p$ & $B_3$ & $p$ & $r$ & $p$ & $\chi^2$ & $p$ \\
        \midrule
        Dreams, 2020 & 0.00 & .467 & \textbf{0.13} & \textbf{.035} & 0.00 & .620 & \textbf{.46} & \textbf{.024} & \textbf{4.48} & \textbf{.034} \\
        Dreams, 2019 & 0.00 & .936 & 0.05 & .538 & 0.00 & .944 & $\oslash$ & $\oslash$ & 0.17 & .679 \\
        Non-dreams, 2020 & 0.00 & .958 & -0.02 & .643 & 0.00 & .943 & .02 & .920 & 0.00 & .969 \\
        \bottomrule
    \end{tabular}
\end{table*}

\section*{Results}

We tracked various forms of dysphoric dream content posted on \textit{r/Dreams} -- a subreddit~\supercite{proferes2021studying} dedicated to dream sharing -- in relation to major events of the first COVID-19 wave (Figure~\ref{fig:summary}A). Across $N = 1,496$ dreams (Figure~\ref{fig:summary}B), we observed a significant increase in anxious dreaming after the World Health Organization (WHO) declared COVID-19 a global pandemic on March 11\textsuperscript{th}, 2020 (interrupted time series analysis~\supercite{bernal2017interrupted}; $B = 0.13$, $CI = $[$0.01$, $0.26$], $p = .035$; Figure~\ref{fig:summary}C). This effect was not present in additional control analyses (Table~\ref{tab:statistics}). We also observed a higher frequency of nightmares in the month following the WHO's announcement than the preceding month (chi-squared test; $\chi^2$ ($1$, $N = 1,494$) $ = 4.3$, $p = .038$; Figure~\ref{fig:summary}D), but not in control analyses (Table~\ref{tab:statistics}). These findings provide strong ecological evidence that the initial COVID-19 global pandemic declaration had a significant impact on dysphoric dreaming and that social media is sensitive to event-driven fluctuations of population dreaming.

\par
To better understand temporal dynamics and the role of media coverage, we examined the relationship between anxious dreaming and COVID-19-related news exposure. Weekly change in COVID-19-related \textit{r/news} posts was predictive of the following week's~\supercite{picard2022targeted} change in anxious dreaming ($r_{22} = 0.45$, $CI =$ [$0.06$, $0.72$], $p = .027$; Figure~\ref{fig:summary}E; see Table~\ref{tab:statistics} for control analyses). This result offers a more fine-grained temporal understanding of COVID-19's impact on dysphoric dreaming and suggests that media exposure might have driven dysphoric dreaming~\supercite{holman2014media}.

\section*{Discussion}

Dysphoric dreaming is only a small part of multidimensional sleep health, but the approach applied in this study might also be effective at tracking other sleep concerns. Dreams are commonly reported with language and thus readily observable on social media~\supercite{elce2021language}, though recent work suggests that Twitter activity can be indicative of sleep loss~\supercite{leypunskiy2018geographically,linnell2021sleep} and more general sleep problems~\supercite{mciver2015characterizing}. Our findings complement these approaches, and further developments will likely reveal that other aspects of sleep health are detectable in social media.

\par
This study has several limitations. First, Reddit users are not necessarily representative of the general population and thus limit our ability to generalize findings. Second, our measure of anxious dreaming may not capture all aspects of dysphoric dreaming. Third, without geographic data attached to Reddit posts, we were unable to use local COVID-19 conditions as a more accurate predictor of changes in dysphoric dreaming.

\par
Currently, tracking sleep health at the population level is a large investment. Traditional survey methods require significant cost and labor to distribute, restricting data collection to discrete time points (e.g., annually). To overcome these limitations, we showed that the language of social media is sensitive to detecting dynamic changes and can be used to estimate population-level dysphoric dream content over time. This low-cost method of digital sleep-health surveillance might offer sleep medicine the opportunity to gain better insights into population sleep patterns and ultimately provide a regional monitoring system for targeted preemptive care.

\section*{Methods}

\textbf{Data source.} Data was extracted from Reddit using the PSAW Python package to interface with Pushshift. Posts from the \textit{r/Dreams} subreddit with ``Short Dream,'' ``Medium Dream,'' or ``Long Dream'' flair were used for dysphoric dreaming measures. Posts from the \textit{r/news} subreddit were used for the COVID-19-related news measure. This research was identified as exempt from review by the Northwestern University Institutional Review Board due to the public nature of this dataset.

\bigskip
\noindent
\textbf{Event selection.} The complete COVID-19 pandemic timeline spans multiple years and includes many milestone events, but for the purposes of this study we required a single event to look for pre/post changes. Based on the authority of the WHO, we used the WHO's declaration of COVID-19 as a global pandemic on March 11\textsuperscript{th}, 2020 for this event Based on prior research suggesting that cultural impacts of an unexpected crisis fade after about 1-2 months~\supercite{pennebaker1993individuals}, we chose a 30-day post-event window (results were broadly consistent at 60-day window).

\bigskip
\noindent
\textbf{Language measures.} Anxious dreaming of each \textit{r/Dreams} post was calculated as the percentage of words belonging to the LIWC-22 anxiety dictionary~\supercite{boyd2022development} (e.g., worry, fear, afraid, nervous). Each \textit{r/Dreams} post was identified as a nightmare if its title contained a non-zero percentage of words belonging to a custom dictionary that included all words beginning with the stem ``nightmar'' (e.g., nightmare, nightmarish). Each \textit{r/news} post was identified as COVID-19-related if its title contained a non-zero percentage of words belonging to a custom COVID-19 dictionary (e.g., covid, virus, mask). All percentages were calculated using the LIWC-22 app~\supercite{boyd2022development}.

\bigskip
\noindent
\textbf{Statistical analyses.} We used an interrupted time series analysis~\supercite{bernal2017interrupted} to test the impact of COVID-19 on anxious dreaming. Anxious dreaming was averaged across all posts for each day surrounding the WHO's declaration of COVID-19 as a global pandemic (March 11\textsuperscript{th}, 2020 $\pm$ 30 days). One day was subtracted from each date of anxious dreaming to account for dreams being reported the day after they occur. A segmented linear regression model was constructed with anxious dreaming as the dependent variable as follows: $Y_t = B_0 + B_1 T + B_2 X_t + B_3 T X_t$, where $T$ is the time elapsed since the first included day and $X_t$ is a dummy variable indicating whether the time was before or after March 11\textsuperscript{th}. Thus, $B_0$ represents the baseline level of anxious dreaming at $T=0$, $B_1$ represents the change in anxious dreaming per-day across the entire time period, $B_2$ represents the level change in anxious dreaming after March 11\textsuperscript{th}, and $B_3$ represents the slope change after March 11th. Model residuals passed visual inspection for autocorrelation, as well as the Durbin-Watson ($d = 2.3$), Breusch-Godfrey ($\chi^2(2) = 3.7, p = .155$), Ljung-Box ($Q(10) = 13.7, p = .186$) tests.

\par
The same March 11\textsuperscript{th} $\pm$ 30-day time window used here was used to compare nightmare frequency before and after March 11\textsuperscript{th}. However, because nightmare was a binary measure (see \textit{Language measures}), post frequencies before and after March 11\textsuperscript{th} were submitted to a 2x2 chi-squared independence test to evaluate change in nightmare frequency.

\par
To test the relationship between COVID-19-related news and anxious dreaming, a larger post-March 11\textsuperscript{th} window (March 12\textsuperscript{th} -- September 1\textsuperscript{st}, 2020) was selected to avoid outliers in COVID-19-related news changes driven by large spikes in the early warning stages ($N_{\text{r/Dreams}} = 5,059$, $N_\text{r/news} = 277,295$). The rate of COVID-19-related news posts and average anxious dreaming were averaged within each week, and weekly anxious dreaming rates were shifted forward $1$ week to account for expected delays in dream incorporation~\supercite{picard2022targeted}. Raw values were converted to percent change between consecutive weeks to account for time series autocorrelation and submitted to a Spearman correlation ($N_\text{weeks} = 24$) because of its lack of normality assumptions. After converting to percent change, residuals passed visual inspection for autocorrelation, as well as the Durbin-Watson ($d = 2.2$), Breusch-Godfrey ($\chi^2(2) = 3.6, p = .169$), Ljung-Box ($Q(10) = 5.8, p = .829$) tests.

\par
All analyses were repeated using the same dates in 2019 and posts from non-dream posts (unflaired posts from \textit{r/Dreams}) as controls for seasonality and general changes in language use, respectively. The interrupted time series and chi-squared tests were repeated at a longer post-event window size (2 months) to control for window selection.

\par
All tests were two-sided. Analyses were implemented in the \texttt{statsmodels} (interrupted time series) and \texttt{Pingouin}~\supercite{vallat2018pingouin} (chi-squared test and correlation) Python packages.

\bigskip
\printbibliography
\end{document}